\section{Component Contributions}
\label{sec:results-ablations}
\label{app:supporting-components}

These component analyses support RQ2 by testing which parts of the graph block
contribute to cross-file detection. The rubric's regeneration variance
(Section~\ref{sec:repro}) makes it too coarse to separate sub-components, so
the analyses use the detection oracle instead.

The deployed graph arm supplies two kinds of dependency evidence
simultaneously: a file-level dependency list and resolved call-graph edges from
the code property graph. Table~\ref{tab:exp2-features} isolates each.

\begin{table}[htbp]
    \centering
    \caption{Detection on structural bands with one kind of dependency
        evidence at a time. Reviews use gpt-4o at $T=0.0$; rates are majority
        votes from five judges at $T=0.0$. The $p$~values are exact McNemar,
        Holm-corrected.}
    \label{tab:exp2-features}
    \begin{tabular}{llcc}
        \toprule
        Arm & Evidence supplied & Detected & Rate \\
        \midrule
        \baseline{}        & none                        & 0/28  & 0\,\% \\
        \texttt{deps only} & file-level dependency list   & 16/28 & 57\,\% \\
        \texttt{edges only}& resolved call edges         & 17/28 & 61\,\% \\
        \kg{}              & both                        & 18/28 & 64\,\% \\
        \midrule
        \multicolumn{2}{l}{Contrast} & $\Delta$ & $p$ (Holm) \\
        \midrule
        \multicolumn{2}{l}{\texttt{deps only} $-$ \baseline{}}  & $+57$pp & $< 0.001$ \\
        \multicolumn{2}{l}{\texttt{edges only} $-$ \baseline{}} & $+61$pp & $< 0.001$ \\
        \bottomrule
    \end{tabular}
\end{table}

Either component alone lifts detection from zero to over half, and both
survive correction. The dependency list recovers sixteen of the eighteen
combined-arm detections, while resolved call edges recover seventeen. Neither
component is significantly below the combined arm at $n=28$. The result
therefore supports partial substitutability, not the necessity or superiority
of resolved edges: the causal claim is that useful structural pointers enable
cross-file detection, while this sample cannot identify one representation as
uniquely responsible.

A pre-registered, Grafana/TypeScript-only test band, evaluated by the earlier
three-judge panel, addressed the test component. On 28 injections where a
renamed symbol provably breaks test files, the filename-convention test
matcher of Section~\ref{sec:approach-kg-render} named a broken test in 24 of 28 cases (86\,\%), compared with 27/28 (96\,\%)
when given the ground truth; the difference is not significant ($p = 0.38$).
Without the test
section, neither \baseline{} nor \kg{} named a broken test in any case. This
shows that the test block contributes detections the dependency block alone
does not.

\subsection{Context Relevance}
\label{sec:results-context-relevance}

The component ablation shows that dependency edges help, but does not rule
out a prompt-length confound: the reviewer might try harder whenever it sees
\emph{any} dependency block, regardless of content. A post-hoc control tests
this possibility.
Table~\ref{tab:exp2-signal-noise} tests this by holding the number of
context items constant and varying only their relevance.
% Source: experiments/2026-09-19_joern_replace_grep/context_ablation/RESULTS_SIGNAL_NOISE_V3.md
% Design follows Yang et al. (EMNLP 2025) GSM-DC matched-distractor methodology.

\begin{table}[htbp]
    \centering
    \caption{Detection under matched-quantity contexts that differ only in
        relevance. \texttt{noise} draws the same number of edges from
        unrelated symbols in the same repository; \texttt{relevant} filters
        the deployed edges to those pointing at true dependents.
        All arms use gpt-4o at $T=0.0$ and the earlier three-judge panel.}
    \label{tab:exp2-signal-noise}
    \begin{tabular}{llcc}
        \toprule
        Arm & Context relevance & Detected & Rate \\
        \midrule
        \baseline{}             & none             & 0/28  & 0\,\% \\
        \texttt{noise}          & wrong edges, matched count & 1/28  & 4\,\% \\
        \kg{}                   & deployed mix     & 18/28 & 64\,\% \\
        \texttt{relevant only}  & true-dependent edges only & 25/28 & 89\,\% \\
        \texttt{idealised}      & ground-truth resolver & 26/28 & 93\,\% \\
        \bottomrule
    \end{tabular}
\end{table}

The noise arm differs little from \baseline{} ($1/28$ versus $0/28$):
matched-volume edges drawn from unrelated symbols provide no comparable
detection gain. Filtering the deployed edges to only those pointing at true
dependents raises detection from 18 to 25 of 28. This post-hoc staircase is
consistent with the context-rot finding that
irrelevant tokens degrade output quality even when the answer is present in
the input~\cite{hong2025contextrot, gupta2025contextlength}.